% Options for packages loaded elsewhere
% Options for packages loaded elsewhere
\PassOptionsToPackage{unicode}{hyperref}
\PassOptionsToPackage{hyphens}{url}
\PassOptionsToPackage{dvipsnames,svgnames,x11names}{xcolor}
%
\documentclass[
  english,
  letterpaper,
  DIV=11,
  numbers=noendperiod]{scrartcl}
\usepackage{xcolor}
\usepackage{amsmath,amssymb}
\setcounter{secnumdepth}{-\maxdimen} % remove section numbering
\usepackage{iftex}
\ifPDFTeX
  \usepackage[T1]{fontenc}
  \usepackage[utf8]{inputenc}
  \usepackage{textcomp} % provide euro and other symbols
\else % if luatex or xetex
  \usepackage{unicode-math} % this also loads fontspec
  \defaultfontfeatures{Scale=MatchLowercase}
  \defaultfontfeatures[\rmfamily]{Ligatures=TeX,Scale=1}
\fi
\usepackage{lmodern}
\ifPDFTeX\else
  % xetex/luatex font selection
\fi
% Use upquote if available, for straight quotes in verbatim environments
\IfFileExists{upquote.sty}{\usepackage{upquote}}{}
\IfFileExists{microtype.sty}{% use microtype if available
  \usepackage[]{microtype}
  \UseMicrotypeSet[protrusion]{basicmath} % disable protrusion for tt fonts
}{}
\makeatletter
\@ifundefined{KOMAClassName}{% if non-KOMA class
  \IfFileExists{parskip.sty}{%
    \usepackage{parskip}
  }{% else
    \setlength{\parindent}{0pt}
    \setlength{\parskip}{6pt plus 2pt minus 1pt}}
}{% if KOMA class
  \KOMAoptions{parskip=half}}
\makeatother
% Make \paragraph and \subparagraph free-standing
\makeatletter
\ifx\paragraph\undefined\else
  \let\oldparagraph\paragraph
  \renewcommand{\paragraph}{
    \@ifstar
      \xxxParagraphStar
      \xxxParagraphNoStar
  }
  \newcommand{\xxxParagraphStar}[1]{\oldparagraph*{#1}\mbox{}}
  \newcommand{\xxxParagraphNoStar}[1]{\oldparagraph{#1}\mbox{}}
\fi
\ifx\subparagraph\undefined\else
  \let\oldsubparagraph\subparagraph
  \renewcommand{\subparagraph}{
    \@ifstar
      \xxxSubParagraphStar
      \xxxSubParagraphNoStar
  }
  \newcommand{\xxxSubParagraphStar}[1]{\oldsubparagraph*{#1}\mbox{}}
  \newcommand{\xxxSubParagraphNoStar}[1]{\oldsubparagraph{#1}\mbox{}}
\fi
\makeatother


\usepackage{longtable,booktabs,array}
\usepackage{calc} % for calculating minipage widths
% Correct order of tables after \paragraph or \subparagraph
\usepackage{etoolbox}
\makeatletter
\patchcmd\longtable{\par}{\if@noskipsec\mbox{}\fi\par}{}{}
\makeatother
% Allow footnotes in longtable head/foot
\IfFileExists{footnotehyper.sty}{\usepackage{footnotehyper}}{\usepackage{footnote}}
\makesavenoteenv{longtable}
\usepackage{graphicx}
\makeatletter
\newsavebox\pandoc@box
\newcommand*\pandocbounded[1]{% scales image to fit in text height/width
  \sbox\pandoc@box{#1}%
  \Gscale@div\@tempa{\textheight}{\dimexpr\ht\pandoc@box+\dp\pandoc@box\relax}%
  \Gscale@div\@tempb{\linewidth}{\wd\pandoc@box}%
  \ifdim\@tempb\p@<\@tempa\p@\let\@tempa\@tempb\fi% select the smaller of both
  \ifdim\@tempa\p@<\p@\scalebox{\@tempa}{\usebox\pandoc@box}%
  \else\usebox{\pandoc@box}%
  \fi%
}
% Set default figure placement to htbp
\def\fps@figure{htbp}
\makeatother


% definitions for citeproc citations
\NewDocumentCommand\citeproctext{}{}
\NewDocumentCommand\citeproc{mm}{%
  \begingroup\def\citeproctext{#2}\cite{#1}\endgroup}
\makeatletter
 % allow citations to break across lines
 \let\@cite@ofmt\@firstofone
 % avoid brackets around text for \cite:
 \def\@biblabel#1{}
 \def\@cite#1#2{{#1\if@tempswa , #2\fi}}
\makeatother
\newlength{\cslhangindent}
\setlength{\cslhangindent}{1.5em}
\newlength{\csllabelwidth}
\setlength{\csllabelwidth}{3em}
\newenvironment{CSLReferences}[2] % #1 hanging-indent, #2 entry-spacing
 {\begin{list}{}{%
  \setlength{\itemindent}{0pt}
  \setlength{\leftmargin}{0pt}
  \setlength{\parsep}{0pt}
  % turn on hanging indent if param 1 is 1
  \ifodd #1
   \setlength{\leftmargin}{\cslhangindent}
   \setlength{\itemindent}{-1\cslhangindent}
  \fi
  % set entry spacing
  \setlength{\itemsep}{#2\baselineskip}}}
 {\end{list}}
\usepackage{calc}
\newcommand{\CSLBlock}[1]{\hfill\break\parbox[t]{\linewidth}{\strut\ignorespaces#1\strut}}
\newcommand{\CSLLeftMargin}[1]{\parbox[t]{\csllabelwidth}{\strut#1\strut}}
\newcommand{\CSLRightInline}[1]{\parbox[t]{\linewidth - \csllabelwidth}{\strut#1\strut}}
\newcommand{\CSLIndent}[1]{\hspace{\cslhangindent}#1}

\ifLuaTeX
\usepackage[bidi=basic]{babel}
\else
\usepackage[bidi=default]{babel}
\fi
% get rid of language-specific shorthands (see #6817):
\let\LanguageShortHands\languageshorthands
\def\languageshorthands#1{}
\ifLuaTeX
  \usepackage[english]{selnolig} % disable illegal ligatures
\fi


\setlength{\emergencystretch}{3em} % prevent overfull lines

\providecommand{\tightlist}{%
  \setlength{\itemsep}{0pt}\setlength{\parskip}{0pt}}



 


\usepackage{fontspec}
\usepackage{multirow}
\usepackage{multicol}
\usepackage{colortbl}
\usepackage[normalem]{ulem}
\usepackage{hhline}
\newlength\Oldarrayrulewidth
\newlength\Oldtabcolsep
\usepackage{longtable}
\usepackage{array}
\usepackage{hyperref}
\usepackage{float}
\usepackage{wrapfig}
\KOMAoption{captions}{tableheading}
\makeatletter
\@ifpackageloaded{caption}{}{\usepackage{caption}}
\AtBeginDocument{%
\ifdefined\contentsname
  \renewcommand*\contentsname{Table of contents}
\else
  \newcommand\contentsname{Table of contents}
\fi
\ifdefined\listfigurename
  \renewcommand*\listfigurename{List of Figures}
\else
  \newcommand\listfigurename{List of Figures}
\fi
\ifdefined\listtablename
  \renewcommand*\listtablename{List of Tables}
\else
  \newcommand\listtablename{List of Tables}
\fi
\ifdefined\figurename
  \renewcommand*\figurename{Figure}
\else
  \newcommand\figurename{Figure}
\fi
\ifdefined\tablename
  \renewcommand*\tablename{Table}
\else
  \newcommand\tablename{Table}
\fi
}
\@ifpackageloaded{float}{}{\usepackage{float}}
\floatstyle{ruled}
\@ifundefined{c@chapter}{\newfloat{codelisting}{h}{lop}}{\newfloat{codelisting}{h}{lop}[chapter]}
\floatname{codelisting}{Listing}
\newcommand*\listoflistings{\listof{codelisting}{List of Listings}}
\makeatother
\makeatletter
\makeatother
\makeatletter
\@ifpackageloaded{caption}{}{\usepackage{caption}}
\@ifpackageloaded{subcaption}{}{\usepackage{subcaption}}
\makeatother
\usepackage{bookmark}
\IfFileExists{xurl.sty}{\usepackage{xurl}}{} % add URL line breaks if available
\urlstyle{same}
\hypersetup{
  pdftitle={Reconstructing individual patient data from published Kaplan-Meier curves: A systematic methodological review and practical guide},
  pdfauthor={Haoyu Ji},
  pdflang={en},
  colorlinks=true,
  linkcolor={blue},
  filecolor={Maroon},
  citecolor={Blue},
  urlcolor={Blue},
  pdfcreator={LaTeX via pandoc}}


\title{Reconstructing individual patient data from published
Kaplan-Meier curves: A systematic methodological review and practical
guide}
\author{Haoyu Ji}
\date{}
\begin{document}
\maketitle


\textbf{Short title:} Reconstruction of survival IPD from Kaplan-Meier
curves\\
\textbf{Affiliations:} {[}Author affiliations{]}\\
\textbf{Corresponding author:} {[}Name and email{]}

\section{Abstract}\label{abstract}

Time-to-event outcomes are central to clinical trials, evidence
synthesis, and health-economic evaluation, but published reports
commonly provide only hazard ratios, selected survival summaries, and
Kaplan-Meier curves. The underlying individual patient data are often
unavailable because of privacy, proprietary, regulatory, and logistical
restrictions. This review identified and classified methods and software
for digitizing published Kaplan-Meier curves and reconstructing
patient-level survival time and event status. Searches of PubMed, Google
Scholar, arXiv, and the Comprehensive R Archive Network identified 698
records; after the prespecified source cap, deduplication, screening,
and full-text assessment, 15 reports were retained and mapped to 11
method or tool entities. The evidence was organized into foundational
algorithms, implementations, core reconstruction methods, upstream
digitization or automation, downstream extensions, and evaluation
studies. IPDfromKM, KMtoIPD, and the CEN-KM component of RESOLVE-IPD
were retained as the three core reconstruction approaches because each
provides distinct reconstruction logic, a runnable implementation, and
patient-level time-to-event output. Their practical suitability depends
mainly on the availability of a risk table, the visibility and overlap
of censoring marks, the presence of an editable vector source, and the
acceptable implementation burden. Other tools serve primarily as
implementations, digitizers, automated interfaces, or extensions for
subgroup and covariate reconstruction. The resulting framework separates
methodological novelty from software packaging and provides a
reproducible decision pathway for selecting, validating, and documenting
a Kaplan-Meier-to-IPD workflow.

\section{1. Introduction}\label{introduction}

Time-to-event outcomes are widely used when the timing of death, disease
progression, recurrence, treatment discontinuation, or another clinical
event is more informative than whether the event occurred alone. In
systematic reviews and meta-analyses, these outcomes are commonly
synthesized with hazard ratios because they incorporate both event
timing and censoring, while Kaplan-Meier (KM) curves provide the
principal graphical summary of the observed survival experience {[}1{]}.
The same information is also central to trial planning,
comparative-effectiveness research, and health-economic models that
require survival over a clinically relevant horizon.

Published summary measures do not preserve all of the information needed
for secondary survival analysis. A hazard ratio is tied to the
proportional-hazards framework and may be difficult to interpret when
treatment effects change over time. Median survival uses only one point
on the survival distribution, and many reports omit alternative
estimands such as fixed-time survival probabilities or restricted mean
survival time (RMST). RMST, defined by the area under the survival curve
up to a prespecified time, can provide a clinically interpretable
treatment contrast without requiring proportional hazards {[}2{]}.
Patient-level time and event data also permit reassessment of
proportional hazards, flexible parametric modeling, estimation of
alternative treatment effects, and analyses that cannot be recovered
from a single published summary {[}3,4{]}.

Individual patient data (IPD) meta-analysis can support harmonized
outcome definitions, consistent analysis across studies, and detailed
assessment of participant- and study-level variation. In practice,
however, original trial IPD are frequently inaccessible. Privacy
protections, proprietary ownership, data-use agreements, regulatory
constraints, and the resources required to curate and share trial data
can all prevent direct access {[}5,6{]}. This creates a recurrent gap
between the analyses that are scientifically desirable and the
information that is publicly available.

Published KM figures preserve part of the missing information. Their
step functions encode event times and changes in survival probability;
risk tables provide intermittent constraints on cohort attrition; total
event counts and plotted censoring symbols may add further information.
Reconstruction methods reverse this reporting process by first
converting the figure into structured coordinates and then assigning
event and censoring times so that the reconstructed patient-level data
reproduce the published curve and accompanying summaries. The iterative
algorithm introduced by Guyot et al.~established the main reconstruction
framework and showed that accuracy improves when numbers at risk or
total event counts are available {[}3{]}. Later implementations improved
preprocessing, usability, numerical stability, treatment of observed
censoring marks, and automation {[}4,5,7,8{]}.

The resulting literature is difficult to interpret as a simple catalogue
of software. A foundational algorithm may appear as author code, a
package, a command in another language, and a later wrapper without
representing four independent statistical methods. Conversely, recent
systems may add a genuinely different reconstruction strategy, automate
only the digitization stage, or target a downstream problem such as
complementary-subgroup recovery, covariate generation, or uncertain
subgroup labels {[}9--12{]}. Validation also differs across studies,
ranging from known simulated IPD to agreement with published summary
measures, making published accuracy statistics unsuitable for an
unqualified ranking.

This study therefore conducted a systematic methodological review of KM
digitization, KM-to-IPD reconstruction, and closely related extensions.
Section 2 describes the search, eligibility criteria, evidence
classification, and report-to-entity mapping. Section 3 compares the
inputs, reconstruction logic, outputs, validation evidence, and
implementation constraints of the three core reconstruction approaches.
Section 4 converts these findings into a practical guide for selecting a
method, validating the reconstructed data, and preserving a reproducible
analysis record.

\section{2. Systematic review methods and evidence
structure}\label{systematic-review-methods-and-evidence-structure}

\subsection{2.1 Review design}\label{review-design}

This systematic methodological review identified and compared methods,
software implementations, and extensions for extracting information from
published KM curves and reconstructing time-to-event IPD. The review is
reported in accordance with PRISMA 2020, with the completed checklist
provided as S1 Checklist {[}13{]}. Because the included evidence did not
estimate a common intervention effect, findings were synthesized
narratively rather than statistically pooled.

\subsection{2.2 Review questions, eligibility criteria, and units of
analysis}\label{review-questions-eligibility-criteria-and-units-of-analysis}

The review addressed three linked questions: which reports proposed or
evaluated KM-specific digitization, KM-to-IPD reconstruction,
censor-identification, or post-reconstruction extension methods; which
independent method or tool entities those records represented; and how
the entities differed in required inputs, reconstruction logic,
censoring treatment, outputs, diagnostics, platform, and
reproducibility.

Eligible records included original reports proposing a KM-specific
digitization, time-to-event reconstruction, censor-identification, or
reconstruction extension method; software packages, commands, web
applications, or research code with sufficient documentation to
reproduce the central workflow; and substantive validation studies.
Application papers that only used reconstructed IPD, generic graph
digitizers without a KM-specific contribution, narrative reviews without
new validation, and records lacking sufficient technical information
were excluded. Publications, packages, web applications, and code
repositories were treated as report or software records. Multiple
records could map to one method or tool entity, while sequential outputs
from the same research group were retained separately when they
introduced materially different reconstruction logic or targets.

\subsection{2.3 Information sources, screening, and data
extraction}\label{information-sources-screening-and-data-extraction}

PubMed, Google Scholar Advanced, arXiv, and the Comprehensive R Archive
Network were searched on 24 May 2026 for records published from 2012
onward. PubMed used three complementary concept groups covering IPD
reconstruction, curve digitization or automation, and method validation.
Google Scholar used four advanced-search combinations with relevance
sorting, while arXiv and the Comprehensive R Archive Network were used
to identify preprints and software records. Complete search strings,
settings, source limits, and returned counts are provided in S1 Table.

The ten logged searches returned 698 records. One Google Scholar query
returned 345 results, of which the first 200 relevance-ranked records
were designated for screening; 145 returned results therefore remained
outside the screening set. After this source-level cap, 553 records were
designated for deduplication and screening. Records were deduplicated
using bibliographic identifiers and normalized titles, screened against
the prespecified eligibility criteria, and assessed in full text. For
eligible records, extraction covered the target task, required figure
and structured inputs, use of risk tables, total events and censor
marks, digitization route, reconstruction algorithm, outputs, diagnostic
measures, platform, implementation status, validation design, and
reported limitations. Report-to-entity mapping was performed after
eligibility assessment so that papers, packages, interfaces, and
wrappers sharing one methodological core were not counted as separate
methods.

\subsection{2.4 Evidence classification and study
selection}\label{evidence-classification-and-study-selection}

The evidence was organized into six functional layers: foundational
algorithms, implementation records, core reconstruction methods,
upstream digitization or end-to-end automation, downstream extensions,
and evaluation or contextual evidence. Core reconstruction was defined
as the conversion of structured KM-derived information into
patient-level records containing at least survival time and event
status. This classification separated methodological novelty from
software packaging and defined the analytic role of each included entity
(Table 1).

\global\setlength{\Oldarrayrulewidth}{\arrayrulewidth}

\global\setlength{\Oldtabcolsep}{\tabcolsep}

\setlength{\tabcolsep}{2pt}

\renewcommand*{\arraystretch}{1.5}



\providecommand{\ascline}[3]{\noalign{\global\arrayrulewidth #1}\arrayrulecolor[HTML]{#2}\cline{#3}}

\begin{longtable}[c]{|p{1.15in}|p{1.35in}|p{2.05in}|p{1.75in}}

\caption{\label{tbl-evidence-layers}Evidence layers and analytic
disposition}

\tabularnewline

\hhline{>{\arrayrulecolor[HTML]{666666}\global\arrayrulewidth=1.5pt}->{\arrayrulecolor[HTML]{666666}\global\arrayrulewidth=1.5pt}->{\arrayrulecolor[HTML]{666666}\global\arrayrulewidth=1.5pt}->{\arrayrulecolor[HTML]{666666}\global\arrayrulewidth=1.5pt}-}

\multicolumn{1}{>{\cellcolor[HTML]{E6E6E6}\raggedright}p{\dimexpr 1.15in+0\tabcolsep}}{\advance\leftskip by 4.0pt\relax \advance\rightskip by 4.0pt\relax \textcolor[HTML]{000000}{\fontsize{9}{9}\selectfont{{\fontspec{Times New Roman} \textbf{Layer}}}}} & \multicolumn{1}{>{\cellcolor[HTML]{E6E6E6}\raggedright}p{\dimexpr 1.35in+0\tabcolsep}}{\advance\leftskip by 4.0pt\relax \advance\rightskip by 4.0pt\relax \textcolor[HTML]{000000}{\fontsize{9}{9}\selectfont{{\fontspec{Times New Roman} \textbf{Representative\ entities}}}}} & \multicolumn{1}{>{\cellcolor[HTML]{E6E6E6}\raggedright}p{\dimexpr 2.05in+0\tabcolsep}}{\advance\leftskip by 4.0pt\relax \advance\rightskip by 4.0pt\relax \textcolor[HTML]{000000}{\fontsize{9}{9}\selectfont{{\fontspec{Times New Roman} \textbf{Defining\ contribution}}}}} & \multicolumn{1}{>{\cellcolor[HTML]{E6E6E6}\raggedright}p{\dimexpr 1.75in+0\tabcolsep}}{\advance\leftskip by 4.0pt\relax \advance\rightskip by 4.0pt\relax \textcolor[HTML]{000000}{\fontsize{9}{9}\selectfont{{\fontspec{Times New Roman} \textbf{Treatment\ in\ the\ review}}}}} \\

\noalign{\global\arrayrulewidth 0pt}\arrayrulecolor[HTML]{000000}

\hhline{>{\arrayrulecolor[HTML]{666666}\global\arrayrulewidth=1.5pt}->{\arrayrulecolor[HTML]{666666}\global\arrayrulewidth=1.5pt}->{\arrayrulecolor[HTML]{666666}\global\arrayrulewidth=1.5pt}->{\arrayrulecolor[HTML]{666666}\global\arrayrulewidth=1.5pt}-}\endfirsthead 

\hhline{>{\arrayrulecolor[HTML]{666666}\global\arrayrulewidth=1.5pt}->{\arrayrulecolor[HTML]{666666}\global\arrayrulewidth=1.5pt}->{\arrayrulecolor[HTML]{666666}\global\arrayrulewidth=1.5pt}->{\arrayrulecolor[HTML]{666666}\global\arrayrulewidth=1.5pt}-}

\multicolumn{1}{>{\cellcolor[HTML]{E6E6E6}\raggedright}p{\dimexpr 1.15in+0\tabcolsep}}{\advance\leftskip by 4.0pt\relax \advance\rightskip by 4.0pt\relax \textcolor[HTML]{000000}{\fontsize{9}{9}\selectfont{{\fontspec{Times New Roman} \textbf{Layer}}}}} & \multicolumn{1}{>{\cellcolor[HTML]{E6E6E6}\raggedright}p{\dimexpr 1.35in+0\tabcolsep}}{\advance\leftskip by 4.0pt\relax \advance\rightskip by 4.0pt\relax \textcolor[HTML]{000000}{\fontsize{9}{9}\selectfont{{\fontspec{Times New Roman} \textbf{Representative\ entities}}}}} & \multicolumn{1}{>{\cellcolor[HTML]{E6E6E6}\raggedright}p{\dimexpr 2.05in+0\tabcolsep}}{\advance\leftskip by 4.0pt\relax \advance\rightskip by 4.0pt\relax \textcolor[HTML]{000000}{\fontsize{9}{9}\selectfont{{\fontspec{Times New Roman} \textbf{Defining\ contribution}}}}} & \multicolumn{1}{>{\cellcolor[HTML]{E6E6E6}\raggedright}p{\dimexpr 1.75in+0\tabcolsep}}{\advance\leftskip by 4.0pt\relax \advance\rightskip by 4.0pt\relax \textcolor[HTML]{000000}{\fontsize{9}{9}\selectfont{{\fontspec{Times New Roman} \textbf{Treatment\ in\ the\ review}}}}} \\

\noalign{\global\arrayrulewidth 0pt}\arrayrulecolor[HTML]{000000}

\hhline{>{\arrayrulecolor[HTML]{666666}\global\arrayrulewidth=1.5pt}->{\arrayrulecolor[HTML]{666666}\global\arrayrulewidth=1.5pt}->{\arrayrulecolor[HTML]{666666}\global\arrayrulewidth=1.5pt}->{\arrayrulecolor[HTML]{666666}\global\arrayrulewidth=1.5pt}-}\endhead



\multicolumn{1}{>{\raggedright}p{\dimexpr 1.15in+0\tabcolsep}}{\advance\leftskip by 4.0pt\relax \advance\rightskip by 4.0pt\relax \textcolor[HTML]{000000}{\fontsize{9}{9}\selectfont{{\fontspec{Times New Roman} Foundational\ algorithm}}}} & \multicolumn{1}{>{\raggedright}p{\dimexpr 1.35in+0\tabcolsep}}{\advance\leftskip by 4.0pt\relax \advance\rightskip by 4.0pt\relax \textcolor[HTML]{000000}{\fontsize{9}{9}\selectfont{{\fontspec{Times New Roman} Guyot/iKM}}}} & \multicolumn{1}{>{\raggedright}p{\dimexpr 2.05in+0\tabcolsep}}{\advance\leftskip by 4.0pt\relax \advance\rightskip by 4.0pt\relax \textcolor[HTML]{000000}{\fontsize{9}{9}\selectfont{{\fontspec{Times New Roman} Original\ iterative\ reconstruction\ logic\ using\ digitized\ Kaplan-Meier\ coordinates\ and\ reported\ risk-set\ information.}}}} & \multicolumn{1}{>{\raggedright}p{\dimexpr 1.75in+0\tabcolsep}}{\advance\leftskip by 4.0pt\relax \advance\rightskip by 4.0pt\relax \textcolor[HTML]{000000}{\fontsize{9}{9}\selectfont{{\fontspec{Times New Roman} Retained\ as\ the\ methodological\ baseline;\ not\ presented\ as\ a\ current\ standalone\ package.}}}} \\

\noalign{\global\arrayrulewidth 0pt}\arrayrulecolor[HTML]{000000}





\multicolumn{1}{>{\raggedright}p{\dimexpr 1.15in+0\tabcolsep}}{\advance\leftskip by 4.0pt\relax \advance\rightskip by 4.0pt\relax \textcolor[HTML]{000000}{\fontsize{9}{9}\selectfont{{\fontspec{Times New Roman} Implementation\ layer}}}} & \multicolumn{1}{>{\raggedright}p{\dimexpr 1.35in+0\tabcolsep}}{\advance\leftskip by 4.0pt\relax \advance\rightskip by 4.0pt\relax \textcolor[HTML]{000000}{\fontsize{9}{9}\selectfont{{\fontspec{Times New Roman} ipdfc;\ survHE;\ reconstructKM}}}} & \multicolumn{1}{>{\raggedright}p{\dimexpr 2.05in+0\tabcolsep}}{\advance\leftskip by 4.0pt\relax \advance\rightskip by 4.0pt\relax \textcolor[HTML]{000000}{\fontsize{9}{9}\selectfont{{\fontspec{Times New Roman} Operational\ implementations\ or\ workflow\ packaging\ of\ Guyot/iKM\ and\ closely\ related\ reconstruction\ steps.}}}} & \multicolumn{1}{>{\raggedright}p{\dimexpr 1.75in+0\tabcolsep}}{\advance\leftskip by 4.0pt\relax \advance\rightskip by 4.0pt\relax \textcolor[HTML]{000000}{\fontsize{9}{9}\selectfont{{\fontspec{Times New Roman} Described\ narratively\ as\ platform\ and\ workflow\ choices;\ not\ counted\ as\ independent\ method\ families.}}}} \\

\noalign{\global\arrayrulewidth 0pt}\arrayrulecolor[HTML]{000000}





\multicolumn{1}{>{\raggedright}p{\dimexpr 1.15in+0\tabcolsep}}{\advance\leftskip by 4.0pt\relax \advance\rightskip by 4.0pt\relax \textcolor[HTML]{000000}{\fontsize{9}{9}\selectfont{{\fontspec{Times New Roman} Core\ reconstruction}}}} & \multicolumn{1}{>{\raggedright}p{\dimexpr 1.35in+0\tabcolsep}}{\advance\leftskip by 4.0pt\relax \advance\rightskip by 4.0pt\relax \textcolor[HTML]{000000}{\fontsize{9}{9}\selectfont{{\fontspec{Times New Roman} IPDfromKM;\ KMtoIPD;\ CEN-KM/RESOLVE-IPD}}}} & \multicolumn{1}{>{\raggedright}p{\dimexpr 2.05in+0\tabcolsep}}{\advance\leftskip by 4.0pt\relax \advance\rightskip by 4.0pt\relax \textcolor[HTML]{000000}{\fontsize{9}{9}\selectfont{{\fontspec{Times New Roman} Distinct\ arm-level\ reconstruction\ logic\ with\ a\ runnable\ implementation\ and\ patient-level\ time\ and\ event\ output.}}}} & \multicolumn{1}{>{\raggedright}p{\dimexpr 1.75in+0\tabcolsep}}{\advance\leftskip by 4.0pt\relax \advance\rightskip by 4.0pt\relax \textcolor[HTML]{000000}{\fontsize{9}{9}\selectfont{{\fontspec{Times New Roman} Compared\ directly\ in\ Section\ 3\ and\ Table\ 2.}}}} \\

\noalign{\global\arrayrulewidth 0pt}\arrayrulecolor[HTML]{000000}





\multicolumn{1}{>{\raggedright}p{\dimexpr 1.15in+0\tabcolsep}}{\advance\leftskip by 4.0pt\relax \advance\rightskip by 4.0pt\relax \textcolor[HTML]{000000}{\fontsize{9}{9}\selectfont{{\fontspec{Times New Roman} Upstream\ digitization\ and\ automation}}}} & \multicolumn{1}{>{\raggedright}p{\dimexpr 1.35in+0\tabcolsep}}{\advance\leftskip by 4.0pt\relax \advance\rightskip by 4.0pt\relax \textcolor[HTML]{000000}{\fontsize{9}{9}\selectfont{{\fontspec{Times New Roman} SurvdigitizeR;\ KM-GPT}}}} & \multicolumn{1}{>{\raggedright}p{\dimexpr 2.05in+0\tabcolsep}}{\advance\leftskip by 4.0pt\relax \advance\rightskip by 4.0pt\relax \textcolor[HTML]{000000}{\fontsize{9}{9}\selectfont{{\fontspec{Times New Roman} Automated\ coordinate\ extraction\ or\ end-to-end\ image\ processing\ around\ a\ reconstruction\ core.}}}} & \multicolumn{1}{>{\raggedright}p{\dimexpr 1.75in+0\tabcolsep}}{\advance\leftskip by 4.0pt\relax \advance\rightskip by 4.0pt\relax \textcolor[HTML]{000000}{\fontsize{9}{9}\selectfont{{\fontspec{Times New Roman} Included\ as\ workflow\ options,\ but\ not\ automatically\ treated\ as\ new\ statistical\ reconstruction\ methods.}}}} \\

\noalign{\global\arrayrulewidth 0pt}\arrayrulecolor[HTML]{000000}





\multicolumn{1}{>{\raggedright}p{\dimexpr 1.15in+0\tabcolsep}}{\advance\leftskip by 4.0pt\relax \advance\rightskip by 4.0pt\relax \textcolor[HTML]{000000}{\fontsize{9}{9}\selectfont{{\fontspec{Times New Roman} Downstream\ extensions}}}} & \multicolumn{1}{>{\raggedright}p{\dimexpr 1.35in+0\tabcolsep}}{\advance\leftskip by 4.0pt\relax \advance\rightskip by 4.0pt\relax \textcolor[HTML]{000000}{\fontsize{9}{9}\selectfont{{\fontspec{Times New Roman} KMSubtraction;\ SynthIPD;\ MAPLE}}}} & \multicolumn{1}{>{\raggedright}p{\dimexpr 2.05in+0\tabcolsep}}{\advance\leftskip by 4.0pt\relax \advance\rightskip by 4.0pt\relax \textcolor[HTML]{000000}{\fontsize{9}{9}\selectfont{{\fontspec{Times New Roman} Recovery\ of\ complementary\ subgroups,\ generation\ of\ covariate-rich\ synthetic\ data,\ or\ propagation\ of\ subgroup-label\ uncertainty.}}}} & \multicolumn{1}{>{\raggedright}p{\dimexpr 1.75in+0\tabcolsep}}{\advance\leftskip by 4.0pt\relax \advance\rightskip by 4.0pt\relax \textcolor[HTML]{000000}{\fontsize{9}{9}\selectfont{{\fontspec{Times New Roman} Discussed\ separately\ because\ the\ target\ differs\ from\ standard\ arm-level\ time-to-event\ reconstruction.}}}} \\

\noalign{\global\arrayrulewidth 0pt}\arrayrulecolor[HTML]{000000}





\multicolumn{1}{>{\raggedright}p{\dimexpr 1.15in+0\tabcolsep}}{\advance\leftskip by 4.0pt\relax \advance\rightskip by 4.0pt\relax \textcolor[HTML]{000000}{\fontsize{9}{9}\selectfont{{\fontspec{Times New Roman} Evaluation\ and\ context}}}} & \multicolumn{1}{>{\raggedright}p{\dimexpr 1.35in+0\tabcolsep}}{\advance\leftskip by 4.0pt\relax \advance\rightskip by 4.0pt\relax \textcolor[HTML]{000000}{\fontsize{9}{9}\selectfont{{\fontspec{Times New Roman} Kim\ et\ al.\ (2025)}}}} & \multicolumn{1}{>{\raggedright}p{\dimexpr 2.05in+0\tabcolsep}}{\advance\leftskip by 4.0pt\relax \advance\rightskip by 4.0pt\relax \textcolor[HTML]{000000}{\fontsize{9}{9}\selectfont{{\fontspec{Times New Roman} Practical\ validation\ or\ comparison\ of\ existing\ reconstruction\ workflows.}}}} & \multicolumn{1}{>{\raggedright}p{\dimexpr 1.75in+0\tabcolsep}}{\advance\leftskip by 4.0pt\relax \advance\rightskip by 4.0pt\relax \textcolor[HTML]{000000}{\fontsize{9}{9}\selectfont{{\fontspec{Times New Roman} Used\ to\ interpret\ accuracy\ and\ implementation\ evidence;\ not\ counted\ as\ a\ new\ method.}}}} \\

\noalign{\global\arrayrulewidth 0pt}\arrayrulecolor[HTML]{000000}

\hhline{>{\arrayrulecolor[HTML]{666666}\global\arrayrulewidth=1.5pt}->{\arrayrulecolor[HTML]{666666}\global\arrayrulewidth=1.5pt}->{\arrayrulecolor[HTML]{666666}\global\arrayrulewidth=1.5pt}->{\arrayrulecolor[HTML]{666666}\global\arrayrulewidth=1.5pt}-}


\end{longtable}

\arrayrulecolor[HTML]{000000}

\global\setlength{\arrayrulewidth}{\Oldarrayrulewidth}

\global\setlength{\tabcolsep}{\Oldtabcolsep}

\renewcommand*{\arraystretch}{1}

\emph{Note:} Classification is based on functional contribution and
reconstruction target rather than author overlap.

After the source-level cap, 553 records entered deduplication.
Twenty-seven duplicates were removed, leaving 526 records for title and
abstract screening. Of these, 507 were excluded and 19 reports were
assessed in full text. Four reports were excluded because they were
application-only reports (n = 2), were not specific to KM digitization
or reconstruction (n = 1), or lacked sufficient technical documentation
(n = 1). Fifteen reports were included and mapped to 11 method or tool
entities. The study-selection sequence is shown in Fig 1.

{[}FIGURE INSERTION: Fig 1. PRISMA 2020 flow diagram{]}

\textbf{Fig 1. PRISMA 2020 flow diagram for identification, screening,
eligibility assessment, and inclusion.} Report-level counts are shown
through full-text inclusion; mapping to method and tool entities was
performed after eligibility assessment.

\subsection{2.5 Boundary of the core
comparison}\label{boundary-of-the-core-comparison}

The Guyot/iKM algorithm was retained as the methodological baseline,
while ipdfc, survHE, and reconstructKM were treated as implementation
records because they operationalize the same foundational reconstruction
family in different languages or broader workflows {[}3,4,14,15{]}.
SurvdigitizeR was retained as an upstream digitization tool because it
returns curve coordinates rather than patient-level event histories
{[}8{]}. KM-GPT was treated as an end-to-end automation system whose
reconstruction stage remains based on an adapted iKM core {[}10{]}.
KMSubtraction, SynthIPD, and MAPLE address complementary-subgroup
recovery, covariate-rich synthetic data, or subgroup-label uncertainty
and therefore target questions beyond standard arm-level reconstruction
{[}9,11,12{]}. These tools remain relevant to the overall workflow but
were not included in the direct comparison of independent arm-level
reconstruction logics.

\section{3. Functional analysis of reconstruction
methods}\label{functional-analysis-of-reconstruction-methods}

\subsection{3.1 Scope and core methods}\label{scope-and-core-methods}

The functional analysis compared capabilities documented in
peer-reviewed articles, preprints, software records, public
repositories, and web documentation available through 13 August 2026. A
method entered the core comparison when it had a distinct and
explainable reconstruction logic, sufficient technical detail, a
publicly runnable implementation, and a primary output containing
patient-level survival time and event status. IPDfromKM, KMtoIPD, and
CEN-KM within RESOLVE-IPD met these criteria. Guyot/iKM remained the
baseline, while implementations, digitizers, automated interfaces, and
downstream extensions were analyzed in their respective layers. Table 2
summarizes the features that directly determine how the three core
methods can be used.

\global\setlength{\Oldarrayrulewidth}{\arrayrulewidth}

\global\setlength{\Oldtabcolsep}{\tabcolsep}

\setlength{\tabcolsep}{2pt}

\renewcommand*{\arraystretch}{1.5}



\providecommand{\ascline}[3]{\noalign{\global\arrayrulewidth #1}\arrayrulecolor[HTML]{#2}\cline{#3}}

\begin{longtable}[c]{|p{1.05in}|p{1.72in}|p{1.72in}|p{1.81in}}

\caption{\label{tbl-core-methods}Core reconstruction methods: inputs,
logic, outputs, and implementation considerations}

\tabularnewline

\hhline{>{\arrayrulecolor[HTML]{666666}\global\arrayrulewidth=1.5pt}->{\arrayrulecolor[HTML]{666666}\global\arrayrulewidth=1.5pt}->{\arrayrulecolor[HTML]{666666}\global\arrayrulewidth=1.5pt}->{\arrayrulecolor[HTML]{666666}\global\arrayrulewidth=1.5pt}-}

\multicolumn{1}{>{\cellcolor[HTML]{E6E6E6}\raggedright}p{\dimexpr 1.05in+0\tabcolsep}}{\advance\leftskip by 4.0pt\relax \advance\rightskip by 4.0pt\relax \textcolor[HTML]{000000}{\fontsize{9}{9}\selectfont{{\fontspec{Times New Roman} \textbf{Domain}}}}} & \multicolumn{1}{>{\cellcolor[HTML]{E6E6E6}\raggedright}p{\dimexpr 1.72in+0\tabcolsep}}{\advance\leftskip by 4.0pt\relax \advance\rightskip by 4.0pt\relax \textcolor[HTML]{000000}{\fontsize{9}{9}\selectfont{{\fontspec{Times New Roman} \textbf{IPDfromKM}}}}} & \multicolumn{1}{>{\cellcolor[HTML]{E6E6E6}\raggedright}p{\dimexpr 1.72in+0\tabcolsep}}{\advance\leftskip by 4.0pt\relax \advance\rightskip by 4.0pt\relax \textcolor[HTML]{000000}{\fontsize{9}{9}\selectfont{{\fontspec{Times New Roman} \textbf{KMtoIPD}}}}} & \multicolumn{1}{>{\cellcolor[HTML]{E6E6E6}\raggedright}p{\dimexpr 1.81in+0\tabcolsep}}{\advance\leftskip by 4.0pt\relax \advance\rightskip by 4.0pt\relax \textcolor[HTML]{000000}{\fontsize{9}{9}\selectfont{{\fontspec{Times New Roman} \textbf{CEN-KM\ /\ RESOLVE-IPD}}}}} \\

\noalign{\global\arrayrulewidth 0pt}\arrayrulecolor[HTML]{000000}

\hhline{>{\arrayrulecolor[HTML]{666666}\global\arrayrulewidth=1.5pt}->{\arrayrulecolor[HTML]{666666}\global\arrayrulewidth=1.5pt}->{\arrayrulecolor[HTML]{666666}\global\arrayrulewidth=1.5pt}->{\arrayrulecolor[HTML]{666666}\global\arrayrulewidth=1.5pt}-}\endfirsthead 

\hhline{>{\arrayrulecolor[HTML]{666666}\global\arrayrulewidth=1.5pt}->{\arrayrulecolor[HTML]{666666}\global\arrayrulewidth=1.5pt}->{\arrayrulecolor[HTML]{666666}\global\arrayrulewidth=1.5pt}->{\arrayrulecolor[HTML]{666666}\global\arrayrulewidth=1.5pt}-}

\multicolumn{1}{>{\cellcolor[HTML]{E6E6E6}\raggedright}p{\dimexpr 1.05in+0\tabcolsep}}{\advance\leftskip by 4.0pt\relax \advance\rightskip by 4.0pt\relax \textcolor[HTML]{000000}{\fontsize{9}{9}\selectfont{{\fontspec{Times New Roman} \textbf{Domain}}}}} & \multicolumn{1}{>{\cellcolor[HTML]{E6E6E6}\raggedright}p{\dimexpr 1.72in+0\tabcolsep}}{\advance\leftskip by 4.0pt\relax \advance\rightskip by 4.0pt\relax \textcolor[HTML]{000000}{\fontsize{9}{9}\selectfont{{\fontspec{Times New Roman} \textbf{IPDfromKM}}}}} & \multicolumn{1}{>{\cellcolor[HTML]{E6E6E6}\raggedright}p{\dimexpr 1.72in+0\tabcolsep}}{\advance\leftskip by 4.0pt\relax \advance\rightskip by 4.0pt\relax \textcolor[HTML]{000000}{\fontsize{9}{9}\selectfont{{\fontspec{Times New Roman} \textbf{KMtoIPD}}}}} & \multicolumn{1}{>{\cellcolor[HTML]{E6E6E6}\raggedright}p{\dimexpr 1.81in+0\tabcolsep}}{\advance\leftskip by 4.0pt\relax \advance\rightskip by 4.0pt\relax \textcolor[HTML]{000000}{\fontsize{9}{9}\selectfont{{\fontspec{Times New Roman} \textbf{CEN-KM\ /\ RESOLVE-IPD}}}}} \\

\noalign{\global\arrayrulewidth 0pt}\arrayrulecolor[HTML]{000000}

\hhline{>{\arrayrulecolor[HTML]{666666}\global\arrayrulewidth=1.5pt}->{\arrayrulecolor[HTML]{666666}\global\arrayrulewidth=1.5pt}->{\arrayrulecolor[HTML]{666666}\global\arrayrulewidth=1.5pt}->{\arrayrulecolor[HTML]{666666}\global\arrayrulewidth=1.5pt}-}\endhead



\multicolumn{1}{>{\raggedright}p{\dimexpr 1.05in+0\tabcolsep}}{\advance\leftskip by 4.0pt\relax \advance\rightskip by 4.0pt\relax \textcolor[HTML]{000000}{\fontsize{9}{9}\selectfont{{\fontspec{Times New Roman} Primary\ target}}}} & \multicolumn{1}{>{\raggedright}p{\dimexpr 1.72in+0\tabcolsep}}{\advance\leftskip by 4.0pt\relax \advance\rightskip by 4.0pt\relax \textcolor[HTML]{000000}{\fontsize{9}{9}\selectfont{{\fontspec{Times New Roman} Arm-level\ time-to-event\ IPD\ from\ a\ published\ Kaplan-Meier\ curve.}}}} & \multicolumn{1}{>{\raggedright}p{\dimexpr 1.72in+0\tabcolsep}}{\advance\leftskip by 4.0pt\relax \advance\rightskip by 4.0pt\relax \textcolor[HTML]{000000}{\fontsize{9}{9}\selectfont{{\fontspec{Times New Roman} Arm-level\ time-to-event\ IPD\ from\ observed\ survival\ drops\ and,\ when\ available,\ marked\ censoring\ times.}}}} & \multicolumn{1}{>{\raggedright}p{\dimexpr 1.81in+0\tabcolsep}}{\advance\leftskip by 4.0pt\relax \advance\rightskip by 4.0pt\relax \textcolor[HTML]{000000}{\fontsize{9}{9}\selectfont{{\fontspec{Times New Roman} Arm-level\ time-to-event\ IPD\ within\ a\ broader\ vector-extraction\ and\ subgroup-analysis\ framework.}}}} \\

\noalign{\global\arrayrulewidth 0pt}\arrayrulecolor[HTML]{000000}





\multicolumn{1}{>{\raggedright}p{\dimexpr 1.05in+0\tabcolsep}}{\advance\leftskip by 4.0pt\relax \advance\rightskip by 4.0pt\relax \textcolor[HTML]{000000}{\fontsize{9}{9}\selectfont{{\fontspec{Times New Roman} Recommended\ input}}}} & \multicolumn{1}{>{\raggedright}p{\dimexpr 1.72in+0\tabcolsep}}{\advance\leftskip by 4.0pt\relax \advance\rightskip by 4.0pt\relax \textcolor[HTML]{000000}{\fontsize{9}{9}\selectfont{{\fontspec{Times New Roman} Digitized\ time-survival\ coordinates,\ baseline\ cohort\ size,\ and\ interval-specific\ numbers\ at\ risk;\ total\ events\ may\ be\ added.}}}} & \multicolumn{1}{>{\raggedright}p{\dimexpr 1.72in+0\tabcolsep}}{\advance\leftskip by 4.0pt\relax \advance\rightskip by 4.0pt\relax \textcolor[HTML]{000000}{\fontsize{9}{9}\selectfont{{\fontspec{Times New Roman} Total\ sample\ size\ and\ lower\ coordinates\ of\ survival\ drops;\ digitized\ censoring\ times\ are\ optional\ but\ central\ when\ no\ risk\ table\ is\ available.}}}} & \multicolumn{1}{>{\raggedright}p{\dimexpr 1.81in+0\tabcolsep}}{\advance\leftskip by 4.0pt\relax \advance\rightskip by 4.0pt\relax \textcolor[HTML]{000000}{\fontsize{9}{9}\selectfont{{\fontspec{Times New Roman} Sample\ size,\ event/drop\ times,\ survival\ probabilities,\ and\ optional\ censoring\ times;\ VEC-KM\ additionally\ requires\ an\ eligible\ vector\ PDF\ or\ SVG.}}}} \\

\noalign{\global\arrayrulewidth 0pt}\arrayrulecolor[HTML]{000000}





\multicolumn{1}{>{\raggedright}p{\dimexpr 1.05in+0\tabcolsep}}{\advance\leftskip by 4.0pt\relax \advance\rightskip by 4.0pt\relax \textcolor[HTML]{000000}{\fontsize{9}{9}\selectfont{{\fontspec{Times New Roman} Risk\ table}}}} & \multicolumn{1}{>{\raggedright}p{\dimexpr 1.72in+0\tabcolsep}}{\advance\leftskip by 4.0pt\relax \advance\rightskip by 4.0pt\relax \textcolor[HTML]{000000}{\fontsize{9}{9}\selectfont{{\fontspec{Times New Roman} Strongly\ recommended\ because\ interval-specific\ risk\ sets\ constrain\ inferred\ censoring\ counts.}}}} & \multicolumn{1}{>{\raggedright}p{\dimexpr 1.72in+0\tabcolsep}}{\advance\leftskip by 4.0pt\relax \advance\rightskip by 4.0pt\relax \textcolor[HTML]{000000}{\fontsize{9}{9}\selectfont{{\fontspec{Times New Roman} Not\ required.}}}} & \multicolumn{1}{>{\raggedright}p{\dimexpr 1.81in+0\tabcolsep}}{\advance\leftskip by 4.0pt\relax \advance\rightskip by 4.0pt\relax \textcolor[HTML]{000000}{\fontsize{9}{9}\selectfont{{\fontspec{Times New Roman} Optional\ in\ the\ documented\ reconstruction\ core;\ current\ risk-table\ alignment\ remains\ less\ mature\ than\ the\ vector/censor\ route.}}}} \\

\noalign{\global\arrayrulewidth 0pt}\arrayrulecolor[HTML]{000000}





\multicolumn{1}{>{\raggedright}p{\dimexpr 1.05in+0\tabcolsep}}{\advance\leftskip by 4.0pt\relax \advance\rightskip by 4.0pt\relax \textcolor[HTML]{000000}{\fontsize{9}{9}\selectfont{{\fontspec{Times New Roman} Censor\ information}}}} & \multicolumn{1}{>{\raggedright}p{\dimexpr 1.72in+0\tabcolsep}}{\advance\leftskip by 4.0pt\relax \advance\rightskip by 4.0pt\relax \textcolor[HTML]{000000}{\fontsize{9}{9}\selectfont{{\fontspec{Times New Roman} Censoring\ is\ inferred\ within\ risk-table\ intervals\ rather\ than\ treating\ each\ plotted\ mark\ as\ an\ observed\ censoring\ time.}}}} & \multicolumn{1}{>{\raggedright}p{\dimexpr 1.72in+0\tabcolsep}}{\advance\leftskip by 4.0pt\relax \advance\rightskip by 4.0pt\relax \textcolor[HTML]{000000}{\fontsize{9}{9}\selectfont{{\fontspec{Times New Roman} Digitized\ censoring\ times\ can\ be\ inserted\ directly;\ ambiguous\ or\ overlapping\ marks\ remain\ a\ source-figure\ limitation.}}}} & \multicolumn{1}{>{\raggedright}p{\dimexpr 1.81in+0\tabcolsep}}{\advance\leftskip by 4.0pt\relax \advance\rightskip by 4.0pt\relax \textcolor[HTML]{000000}{\fontsize{9}{9}\selectfont{{\fontspec{Times New Roman} Explicit\ and\ duplicate\ censoring\ times\ can\ be\ represented\ while\ candidate\ allocations\ are\ evaluated\ against\ the\ target\ curve.}}}} \\

\noalign{\global\arrayrulewidth 0pt}\arrayrulecolor[HTML]{000000}





\multicolumn{1}{>{\raggedright}p{\dimexpr 1.05in+0\tabcolsep}}{\advance\leftskip by 4.0pt\relax \advance\rightskip by 4.0pt\relax \textcolor[HTML]{000000}{\fontsize{9}{9}\selectfont{{\fontspec{Times New Roman} Digitization\ route}}}} & \multicolumn{1}{>{\raggedright}p{\dimexpr 1.72in+0\tabcolsep}}{\advance\leftskip by 4.0pt\relax \advance\rightskip by 4.0pt\relax \textcolor[HTML]{000000}{\fontsize{9}{9}\selectfont{{\fontspec{Times New Roman} Interactive\ point\ acquisition\ in\ the\ R/Shiny\ workflow\ or\ externally\ digitized\ coordinates.}}}} & \multicolumn{1}{>{\raggedright}p{\dimexpr 1.72in+0\tabcolsep}}{\advance\leftskip by 4.0pt\relax \advance\rightskip by 4.0pt\relax \textcolor[HTML]{000000}{\fontsize{9}{9}\selectfont{{\fontspec{Times New Roman} External\ digitization\ of\ the\ lower\ step\ coordinates\ and\ censor\ marks.}}}} & \multicolumn{1}{>{\raggedright}p{\dimexpr 1.81in+0\tabcolsep}}{\advance\leftskip by 4.0pt\relax \advance\rightskip by 4.0pt\relax \textcolor[HTML]{000000}{\fontsize{9}{9}\selectfont{{\fontspec{Times New Roman} VEC-KM\ extracts\ curves\ and\ censor\ symbols\ from\ vector\ graphics;\ target-panel\ identification\ is\ still\ required.}}}} \\

\noalign{\global\arrayrulewidth 0pt}\arrayrulecolor[HTML]{000000}





\multicolumn{1}{>{\raggedright}p{\dimexpr 1.05in+0\tabcolsep}}{\advance\leftskip by 4.0pt\relax \advance\rightskip by 4.0pt\relax \textcolor[HTML]{000000}{\fontsize{9}{9}\selectfont{{\fontspec{Times New Roman} Reconstruction\ logic}}}} & \multicolumn{1}{>{\raggedright}p{\dimexpr 1.72in+0\tabcolsep}}{\advance\leftskip by 4.0pt\relax \advance\rightskip by 4.0pt\relax \textcolor[HTML]{000000}{\fontsize{9}{9}\selectfont{{\fontspec{Times New Roman} Modified\ iterative\ Guyot/iKM\ algorithm\ with\ automated\ preprocessing\ and\ numerical\ safeguards.}}}} & \multicolumn{1}{>{\raggedright}p{\dimexpr 1.72in+0\tabcolsep}}{\advance\leftskip by 4.0pt\relax \advance\rightskip by 4.0pt\relax \textcolor[HTML]{000000}{\fontsize{9}{9}\selectfont{{\fontspec{Times New Roman} Greedy\ event\ allocation\ at\ observed\ Kaplan-Meier\ drops\ with\ optional\ direct\ use\ of\ censoring\ times.}}}} & \multicolumn{1}{>{\raggedright}p{\dimexpr 1.81in+0\tabcolsep}}{\advance\leftskip by 4.0pt\relax \advance\rightskip by 4.0pt\relax \textcolor[HTML]{000000}{\fontsize{9}{9}\selectfont{{\fontspec{Times New Roman} CEN-KM\ branch-and-compare\ search\ over\ censoring\ configurations,\ followed\ by\ target-curve\ comparison.}}}} \\

\noalign{\global\arrayrulewidth 0pt}\arrayrulecolor[HTML]{000000}





\multicolumn{1}{>{\raggedright}p{\dimexpr 1.05in+0\tabcolsep}}{\advance\leftskip by 4.0pt\relax \advance\rightskip by 4.0pt\relax \textcolor[HTML]{000000}{\fontsize{9}{9}\selectfont{{\fontspec{Times New Roman} Primary\ output}}}} & \multicolumn{1}{>{\raggedright}p{\dimexpr 1.72in+0\tabcolsep}}{\advance\leftskip by 4.0pt\relax \advance\rightskip by 4.0pt\relax \textcolor[HTML]{000000}{\fontsize{9}{9}\selectfont{{\fontspec{Times New Roman} Patient-level\ survival\ time\ and\ event\ status,\ reconstructed\ curve,\ risk-set\ comparison,\ and\ standard\ survival\ summaries.}}}} & \multicolumn{1}{>{\raggedright}p{\dimexpr 1.72in+0\tabcolsep}}{\advance\leftskip by 4.0pt\relax \advance\rightskip by 4.0pt\relax \textcolor[HTML]{000000}{\fontsize{9}{9}\selectfont{{\fontspec{Times New Roman} Patient-level\ survival\ time\ and\ event\ indicator.}}}} & \multicolumn{1}{>{\raggedright}p{\dimexpr 1.81in+0\tabcolsep}}{\advance\leftskip by 4.0pt\relax \advance\rightskip by 4.0pt\relax \textcolor[HTML]{000000}{\fontsize{9}{9}\selectfont{{\fontspec{Times New Roman} Patient-level\ survival\ time\ and\ event\ indicator;\ MAPLE\ can\ add\ ensembles\ of\ subgroup\ labels\ in\ the\ broader\ framework.}}}} \\

\noalign{\global\arrayrulewidth 0pt}\arrayrulecolor[HTML]{000000}





\multicolumn{1}{>{\raggedright}p{\dimexpr 1.05in+0\tabcolsep}}{\advance\leftskip by 4.0pt\relax \advance\rightskip by 4.0pt\relax \textcolor[HTML]{000000}{\fontsize{9}{9}\selectfont{{\fontspec{Times New Roman} Diagnostics}}}} & \multicolumn{1}{>{\raggedright}p{\dimexpr 1.72in+0\tabcolsep}}{\advance\leftskip by 4.0pt\relax \advance\rightskip by 4.0pt\relax \textcolor[HTML]{000000}{\fontsize{9}{9}\selectfont{{\fontspec{Times New Roman} Curve\ and\ risk-set\ comparison\ with\ RMSE,\ MAE,\ maximum\ absolute\ error,\ Kolmogorov-Smirnov\ statistic,\ and\ survival\ summaries.}}}} & \multicolumn{1}{>{\raggedright}p{\dimexpr 1.72in+0\tabcolsep}}{\advance\leftskip by 4.0pt\relax \advance\rightskip by 4.0pt\relax \textcolor[HTML]{000000}{\fontsize{9}{9}\selectfont{{\fontspec{Times New Roman} The\ method\ paper\ evaluates\ curve\ fit,\ hazard\ ratios,\ and\ survival\ quantiles;\ a\ standardized\ package\ report\ is\ not\ built\ in.}}}} & \multicolumn{1}{>{\raggedright}p{\dimexpr 1.81in+0\tabcolsep}}{\advance\leftskip by 4.0pt\relax \advance\rightskip by 4.0pt\relax \textcolor[HTML]{000000}{\fontsize{9}{9}\selectfont{{\fontspec{Times New Roman} Target-curve\ tolerance\ and\ refitting\ checks;\ research-code\ outputs\ require\ explicit\ archiving\ of\ inputs\ and\ settings.}}}} \\

\noalign{\global\arrayrulewidth 0pt}\arrayrulecolor[HTML]{000000}





\multicolumn{1}{>{\raggedright}p{\dimexpr 1.05in+0\tabcolsep}}{\advance\leftskip by 4.0pt\relax \advance\rightskip by 4.0pt\relax \textcolor[HTML]{000000}{\fontsize{9}{9}\selectfont{{\fontspec{Times New Roman} Implementation\ considerations}}}} & \multicolumn{1}{>{\raggedright}p{\dimexpr 1.72in+0\tabcolsep}}{\advance\leftskip by 4.0pt\relax \advance\rightskip by 4.0pt\relax \textcolor[HTML]{000000}{\fontsize{9}{9}\selectfont{{\fontspec{Times New Roman} Peer-reviewed\ R\ package\ and\ Shiny\ workflow;\ well\ suited\ to\ reproducible\ local\ analysis\ when\ risk-table\ information\ is\ available.}}}} & \multicolumn{1}{>{\raggedright}p{\dimexpr 1.72in+0\tabcolsep}}{\advance\leftskip by 4.0pt\relax \advance\rightskip by 4.0pt\relax \textcolor[HTML]{000000}{\fontsize{9}{9}\selectfont{{\fontspec{Times New Roman} Peer-reviewed\ R\ implementation\ with\ a\ compact\ workflow;\ performance\ depends\ on\ reliable\ identification\ of\ drops\ and\ censor\ marks.}}}} & \multicolumn{1}{>{\raggedright}p{\dimexpr 1.81in+0\tabcolsep}}{\advance\leftskip by 4.0pt\relax \advance\rightskip by 4.0pt\relax \textcolor[HTML]{000000}{\fontsize{9}{9}\selectfont{{\fontspec{Times New Roman} Preprint\ and\ research-code\ framework\ with\ vector-specific\ advantages\ and\ greater\ version-control\ and\ reproducibility\ burden.}}}} \\

\noalign{\global\arrayrulewidth 0pt}\arrayrulecolor[HTML]{000000}

\hhline{>{\arrayrulecolor[HTML]{666666}\global\arrayrulewidth=1.5pt}->{\arrayrulecolor[HTML]{666666}\global\arrayrulewidth=1.5pt}->{\arrayrulecolor[HTML]{666666}\global\arrayrulewidth=1.5pt}->{\arrayrulecolor[HTML]{666666}\global\arrayrulewidth=1.5pt}-}


\end{longtable}

\arrayrulecolor[HTML]{000000}

\global\setlength{\arrayrulewidth}{\Oldarrayrulewidth}

\global\setlength{\tabcolsep}{\Oldtabcolsep}

\renewcommand*{\arraystretch}{1}

\emph{Note:} The comparison is documentation-based and does not
represent a new head-to-head execution benchmark.

\subsection{3.2 Comparative characteristics, validation, and practical
limitations}\label{comparative-characteristics-validation-and-practical-limitations}

IPDfromKM combines interactive coordinate acquisition, automated
preprocessing, a modified iKM reconstruction algorithm, and diagnostic
summaries in an R package and Shiny workflow {[}5{]}. Its recommended
inputs are digitized time-survival coordinates, baseline cohort size,
and risk-table times with corresponding numbers at risk; a total event
count may provide an additional constraint. The risk table is central to
the method because interval-specific risk sets constrain the number of
inferred censoring observations. Although the algorithm can run with
baseline cohort size alone, the censoring pattern is then less
identifiable and the reconstruction is more weakly constrained.

KMtoIPD uses the total sample size and the lower coordinates of observed
survival drops to allocate event counts and can insert digitized
censoring times directly into the patient-level sequence {[}7{]}. It
therefore does not require a risk table and is most informative when
censoring marks are clearly visible and can be separated from curve
steps. Its use of plotted censoring information does not remove
ambiguity from the source figure: one mark may conceal multiple
censoring events, marks may overlap, and raster compression or multiple
curves may prevent reliable assignment.

RESOLVE-IPD separates high-fidelity vector extraction from
censoring-informed reconstruction {[}12{]}. VEC-KM extracts curves and
censor symbols from eligible vector PDF or SVG content, while CEN-KM
evaluates candidate censoring configurations against the target curve
through a branch-and-compare procedure. This permits explicit and
duplicate censoring times and differs from both risk-table-constrained
iKM reconstruction and the adapted-iKM pipeline used by KM-GPT. The
method's vector-specific advantages are accompanied by research-code
maturity, target-panel identification, and a greater need to freeze
versions, parameters, and intermediate files.

All three methods return patient-level survival time and event status.
IPDfromKM additionally supplies reconstructed curves, risk-set
comparisons, RMSE, MAE, maximum absolute error, a Kolmogorov-Smirnov
statistic, and common survival summaries {[}5{]}. KMtoIPD returns a
compact patient-level dataset; the accompanying study evaluates curve
agreement, hazard ratios, and survival quantiles, so a consistent
diagnostic report must be constructed around the output {[}7{]}.
RESOLVE-IPD uses target-curve tolerance and refitting checks and can add
ensembles of plausible subgroup labels through MAPLE, although those
label outputs belong to the extension layer rather than the core
arm-level reconstruction. Diagnostic thresholds proposed in individual
papers were treated as method-specific empirical rules rather than
universal acceptance criteria.

Validation evidence differed in target and reference standard.
Guyot/iKM, IPDfromKM, and KMtoIPD used simulated data, original IPD, or
published summaries to evaluate reconstructed curves and survival
estimates {[}3,5,7{]}. SurvdigitizeR evaluated coordinate extraction
rather than the complete patient-level reconstruction {[}8{]}, while the
principal evidence for KM-GPT, SynthIPD, and RESOLVE-IPD remained
author-led preprint validation at the evidence freeze {[}10--12{]}. The
real-world study by Kim et al.~provided useful external evidence for an
existing workflow, but differences in figures, information availability,
sample size, censoring, and evaluation metrics prevented a direct
cross-method ranking {[}6{]}. Simulation evidence was interpreted
according to the aims, data-generating mechanisms, estimands, methods,
and performance measures reported by each study {[}16{]}.

Implementation choices alter the practical workflow without changing the
underlying method family. ipdfc provides a peer-reviewed Stata
implementation of iKM, survHE embeds reconstruction within a
health-economic survival-modeling suite, and reconstructKM supplies R
utilities for formatted curve and risk-table inputs {[}4,14,15{]}. Local
R, Stata, or Python workflows support versioned scripts and offline
archiving, whereas browser and model-assisted systems may introduce
service availability, data-transfer, credential, and model-version
dependencies. These differences should be recorded as reproducibility
characteristics rather than interpreted as evidence of statistical
superiority.

Across methods, the most consequential failures arise before or during
reconstruction: the wrong curve may be matched to a risk table, axes may
be clipped or miscalibrated, censoring symbols may be unobservable, and
low-resolution images may prevent reproducible localization of curve
steps. The response depends on whether the missing information can be
replaced by another reliable constraint. A reconstruction should be
revised or abandoned when source identity, calibration, or validation
cannot be defended (Fig 2).

{[}FIGURE INSERTION: Fig 2 from KM\_IPD\_review\_figures.pptx{]}

\textbf{Fig 2. Common failure signals and response rules for KM-to-IPD
reconstruction.} The compact framework links missing risk information,
ambiguous censoring, poor image quality, and curve-risk mismatch to an
alternative input route, redigitization, sensitivity analysis, or an
explicit stop decision.

\section{4. Practical guide for method selection and
reconstruction}\label{practical-guide-for-method-selection-and-reconstruction}

\subsection{4.1 Define the target
output}\label{define-the-target-output}

Method selection begins with the intended patient-level output. Standard
arm-level reconstruction seeks a dataset containing survival time, event
status, and an arm identifier for each reconstructed patient. A
different workflow is required when the target is an unreported
complementary subgroup: the overall cohort and a known subgroup must
first be reconstructed before KMSubtraction can estimate the remaining
subgroup {[}9{]}. Covariate-rich synthetic data or uncertain subgroup
labels require additional published subgroup summaries and belong to the
SynthIPD or MAPLE extension layer rather than ordinary arm-level
reconstruction {[}11,12{]}.

\subsection{4.2 Match the method to the source figure and available
constraints}\label{match-the-method-to-the-source-figure-and-available-constraints}

For standard arm-level IPD, the first decision is whether an editable
vector figure is available. Eligible vector PDF or SVG content can be
routed through VEC-KM and CEN-KM when precise extraction of curve steps
and censor symbols is important and research-code use is acceptable
{[}12{]}. Raster images, screenshots, and non-editable PDFs require
manual or automated digitization before reconstruction. SurvdigitizeR
can reduce the manual burden of raster coordinate extraction, whereas
KM-GPT offers an end-to-end image-to-IPD route when an external model,
web service, and data transfer are acceptable {[}8,10{]}.

Within a raster or externally digitized workflow, a reliable risk table
favors IPDfromKM or another documented iKM implementation because
interval-specific risk sets directly constrain inferred censoring
{[}3,5{]}. When no risk table is available but survival drops and censor
marks are clearly distinguishable, KMtoIPD can use the visible censoring
information directly {[}7{]}. Overlapping or ambiguous censor marks
favor a vector/CEN-KM route when available; otherwise, alternative
plausible allocations should be examined and the resulting uncertainty
reported. Reconstruction should not proceed when the axes cannot be
calibrated or when the curve cannot be uniquely matched to its legend,
arm, and risk table. The complete decision pathway is summarized in Fig
3.

{[}FIGURE INSERTION: Fig 3 from KM\_IPD\_review\_figures.pptx{]}

\textbf{Fig 3. Method-selection pathway for reconstruction from
published Kaplan-Meier figures.} The pathway starts with the target
output and source format, then uses risk-table availability, censor-mark
quality, vector eligibility, automation requirements, and implementation
constraints to select a reconstruction route.

\subsection{4.3 Apply a common reconstruction and validation
workflow}\label{apply-a-common-reconstruction-and-validation-workflow}

Regardless of the selected tool, the analysis should preserve a common
sequence. The source record should identify the article version, figure
and panel, treatment arm, file type, and date of acquisition. The target
panel and risk table should be isolated, the axes calibrated to the
reported time and survival scales, and curve and censor coordinates
extracted without discarding the original image or raw points. Cleaning
should sort coordinates by time, remove exact duplicates, enforce valid
survival bounds, document any monotonicity correction, and verify that
curve, legend, arm label, baseline sample size, and risk-table times
refer to the same population.

Reconstruction should be run with all available constraints and then
evaluated by re-estimating the KM curve from the patient-level data. At
minimum, validation should include a visual overlay and comparison of
reported numbers at risk when present. Fixed-time survival
probabilities, median survival, hazard ratios, or RMST should also be
compared when the publication provides compatible summaries {[}3,5,6{]}.
Disagreement should trigger review of panel isolation, axis calibration,
digitization, risk-table transcription, censor interpretation, and
method assumptions rather than automatic acceptance of a
complete-looking dataset.

The final archive should contain the source figure, raw and cleaned
coordinates, risk-table and censor inputs, reconstruction parameters,
patient-level output, diagnostic plots and summaries, software and
dependency versions, scripts, and a log of manual changes. For
stochastic, web-based, or model-assisted workflows, the random seed,
model or service version, execution date, and any unavailable external
dependency should also be recorded. This common record allows the
reconstruction to be repeated, audited, or replaced when a
higher-quality source or improved method becomes available.

\section{References}\label{references}

\phantomsection\label{refs}
\begin{CSLReferences}{0}{1}
\bibitem[\citeproctext]{ref-tierney2007}
\CSLLeftMargin{1. }%
\CSLRightInline{Tierney JF, Stewart LA, Ghersi D, Burdett S, Sydes MR.
\href{https://doi.org/10.1186/1745-6215-8-16}{Practical methods for
incorporating summary time-to-event data into meta-analysis}. Trials.
2007;8:16. }

\bibitem[\citeproctext]{ref-royston2013}
\CSLLeftMargin{2. }%
\CSLRightInline{Royston P, Parmar MKB.
\href{https://doi.org/10.1186/1471-2288-13-152}{Restricted mean survival
time: An alternative to the hazard ratio for the design and analysis of
randomized trials with a time-to-event outcome}. BMC Medical Research
Methodology. 2013;13:152. }

\bibitem[\citeproctext]{ref-guyot2012}
\CSLLeftMargin{3. }%
\CSLRightInline{Guyot P, Ades AE, Ouwens MJNM, Welton NJ.
\href{https://doi.org/10.1186/1471-2288-12-9}{Enhanced secondary
analysis of survival data: Reconstructing the data from published
kaplan-meier survival curves}. BMC Medical Research Methodology.
2012;12:9. }

\bibitem[\citeproctext]{ref-wei2017}
\CSLLeftMargin{4. }%
\CSLRightInline{Wei Y, Royston P.
\href{https://doi.org/10.1177/1536867X1801700402}{Reconstructing
time-to-event data from published kaplan-meier curves}. The Stata
Journal. 2017;17(4):786--802. }

\bibitem[\citeproctext]{ref-liu2021}
\CSLLeftMargin{5. }%
\CSLRightInline{Liu N, Zhou Y, Lee JJ.
\href{https://doi.org/10.1186/s12874-021-01308-8}{{IPDfromKM}:
Reconstruct individual patient data from published kaplan-meier survival
curves}. BMC Medical Research Methodology. 2021;21:111. }

\bibitem[\citeproctext]{ref-kim2025}
\CSLLeftMargin{6. }%
\CSLRightInline{Kim J, Chalise P, He J.
\href{https://doi.org/10.1016/j.conctc.2025.101542}{Reconstructing
patient level survival data from published kaplan-meier curves}.
Contemporary Clinical Trials Communications. 2025;47:101542. }

\bibitem[\citeproctext]{ref-rogula2022}
\CSLLeftMargin{7. }%
\CSLRightInline{Rogula B, Lozano-Ortega G, Johnston KM.
\href{https://doi.org/10.1177/23814683221077643}{A method for
reconstructing individual patient data from kaplan-meier survival curves
that incorporate marked censoring times}. MDM Policy \& Practice.
2022;7(1):23814683221077643. }

\bibitem[\citeproctext]{ref-zhang2024}
\CSLLeftMargin{8. }%
\CSLRightInline{Zhang JZ, Rios JD, Pechlivanoglou T, Yang A, Zhang Q,
Deris D, et al.
\href{https://doi.org/10.1186/s12874-024-02273-8}{{SurvdigitizeR}: An
algorithm for automated survival curve digitization}. BMC Medical
Research Methodology. 2024;24:147. }

\bibitem[\citeproctext]{ref-zhao2022}
\CSLLeftMargin{9. }%
\CSLRightInline{Zhao J, Syn NL, Tan BKJ, Yap DWT, Teo CB, Chan YH, et
al. \href{https://doi.org/10.1186/s12874-022-01567-z}{{KMSubtraction}:
Reconstruction of unreported subgroup survival data utilizing published
kaplan-meier survival curves}. BMC Medical Research Methodology.
2022;22:93. }

\bibitem[\citeproctext]{ref-kmgpt2025}
\CSLLeftMargin{10. }%
\CSLRightInline{Zhao Y, Sun H, Ding Y, Xu Y. {KM-GPT}: An automated
pipeline for reconstructing individual patient data from kaplan-meier
plots {[}Internet{]}. 2025. Available from:
\url{https://arxiv.org/abs/2509.18141}}

\bibitem[\citeproctext]{ref-synthipd2026}
\CSLLeftMargin{11. }%
\CSLRightInline{Zhao Z, Ren Z, Zhai G, Hu F, Ma W, Xie E, et al.
{SynthIPD}: Training-free synthetic individual patient data generation
{[}Internet{]}. 2026. Available from:
\url{https://arxiv.org/abs/2509.16466}}

\bibitem[\citeproctext]{ref-resolve2026}
\CSLLeftMargin{12. }%
\CSLRightInline{Lang L, Zhao Y, Gao Q, Xu Y. {RESOLVE-IPD}:
High-fidelity individual patient data reconstruction and
uncertainty-aware subgroup meta-analysis {[}Internet{]}. 2026. Available
from: \url{https://arxiv.org/abs/2511.01785}}

\bibitem[\citeproctext]{ref-page2021}
\CSLLeftMargin{13. }%
\CSLRightInline{Page MJ, McKenzie JE, Bossuyt PM, Boutron I, Hoffmann
TC, Mulrow CD, et al. \href{https://doi.org/10.1136/bmj.n71}{The
{PRISMA} 2020 statement: An updated guideline for reporting systematic
reviews}. BMJ. 2021;372:n71. }

\bibitem[\citeproctext]{ref-baio2020}
\CSLLeftMargin{14. }%
\CSLRightInline{Baio G.
\href{https://doi.org/10.18637/jss.v095.i14}{{survHE}: Survival analysis
for health economic evaluation and cost-effectiveness modeling}. Journal
of Statistical Software. 2020;95(14):1--47. }

\bibitem[\citeproctext]{ref-reconstructkm2025}
\CSLLeftMargin{15. }%
\CSLRightInline{Sun R. {reconstructKM}: Reconstruct individual-level
data from published kaplan-meier plots {[}Internet{]}. 2025. Available
from: \url{https://cran.r-project.org/package=reconstructKM}}

\bibitem[\citeproctext]{ref-morris2019}
\CSLLeftMargin{16. }%
\CSLRightInline{Morris TP, White IR, Crowther MJ.
\href{https://doi.org/10.1002/sim.8086}{Using simulation studies to
evaluate statistical methods}. Statistics in Medicine.
2019;38(11):2074--102. }

\end{CSLReferences}

\section{Supporting information
captions}\label{supporting-information-captions}

\textbf{S1 Checklist. PRISMA 2020 checklist.} Completed reporting
checklist for the systematic methodological review.

\textbf{S1 Table. Full search strategies and source-level counts.} Exact
database and platform queries, search dates, limits, raw hits, and the
Google Scholar source cap.

\textbf{S1 Appendix. Review protocol and amendments.} Review questions,
eligibility criteria, extraction fields, report-to-entity mapping rules,
and amendment history.




\end{document}
